\documentclass[11pt]{article}

\usepackage[margin=1in]{geometry}
\usepackage{amsmath,amssymb,amsthm,mathtools}
\usepackage[T1]{fontenc}
\usepackage{lmodern}
\usepackage{microtype}
\usepackage{hyperref}

\hypersetup{
  colorlinks=true,
  linkcolor=blue,
  urlcolor=blue
}

\newcommand{\RR}{\mathbb{R}}
\newcommand{\norm}[1]{\left\lVert #1 \right\rVert}
\newcommand{\vect}[1]{\mathbf{#1}}

\title{A Geodesic Variant of the Kamada--Kawai Energy}
\author{}
\date{\today}

\begin{document}

\maketitle

\section*{Introduction}

The classical Kamada--Kawai (KK) energy is built from the discrepancy between graph
distances and ambient Euclidean distances between embedded vertices. For an embedding
$X = (\vect{x}_1,\dots,\vect{x}_n)$ with $\vect{x}_i \in \RR^d$, the standard model uses
\[
E_{\mathrm{KK}}(X)
=
\frac{1}{2}\sum_{1 \le i < j \le n}
k_{ij}\bigl(\norm{\vect{x}_i-\vect{x}_j}-\ell_{ij}\bigr)^2,
\]
where
\[
\ell_{ij} = L_0 g_{ij},
\qquad
L_0 = \frac{L}{D},
\qquad
k_{ij} = \frac{K}{g_{ij}^2}.
\]
Here $g_{ij}$ is the graph distance between vertices $i$ and $j$, $D$ is the graph
diameter, $L$ is the drawing scale, and $K$ is a global stiffness constant.

This formulation is natural if the goal is to match graph-theoretic distances with
straight-line distances in the ambient space. However, the graph itself carries an intrinsic
one-dimensional geometry: movement from $i$ to $j$ takes place along graph edges rather
than along the direct chord joining $\vect{x}_i$ and $\vect{x}_j$ in $\RR^d$. Because of this,
the ambient Euclidean distance $\norm{\vect{x}_i-\vect{x}_j}$ can be too crude. Two vertices
may be far apart along the graph but happen to lie close in the drawing, or vice versa.

A more intrinsic alternative is to measure distance \emph{along the embedded graph}. In
other words, once the graph is drawn in $\RR^d$, each edge inherits its Euclidean length,
and the distance between two vertices is measured as the total length of a graph path in the
embedding. This leads to a geodesic version of the KK energy that compares graph distance
with embedding-induced path length rather than with naive straight-line separation.

Classical KK and stress-majorization methods use graph distances as \emph{targets}, but the
realized distance in the optimized functional is still the ambient Euclidean distance between
vertex positions \cite{kamada1989algorithm,gansner2004stress}. Cohen proposed
linear-network distance as an alternative target metric, but again compared that target to
Euclidean inter-node distance rather than to path length measured along the drawn graph
\cite{cohen1997drawing}; see also the explicit discussion in \cite{koren2003axis}. Sparse
stress models are close in spirit to the landmark sparsification proposed here, but they
likewise retain the standard Euclidean discrepancy term \cite{ortmann2017sparse}. A recent
path-aware layout method, Path-Guided Stochastic Gradient Descent, also optimizes Euclidean
layout distances against path-derived targets rather than against embedded path lengths
\cite{heumos2024pangenome}. Outside graph drawing, Isomap is a conceptual cousin in that it
uses graph geodesics in an MDS-style embedding, but those geodesics come from a fixed input
graph rather than from edge lengths induced by the current drawing \cite{tenenbaum2000global}.
In particular, \cite{gansner2004stress} is not a formulation of the present geodesic KK
functional: its stress term is still based on $\norm{\vect{x}_i-\vect{x}_j}$, not on
$h_{ij}(X)$.

\section*{Embedding-Induced Geodesic Distance}

Let $G=(V,E)$ be a connected graph with $V=\{1,\dots,n\}$. For every unordered pair
$i<j$, fix once and for all a graph shortest path
\[
\gamma_{ij} \subseteq E
\]
between $i$ and $j$, where shortest path is understood with respect to the graph metric
underlying $g_{ij}$. In an unweighted graph this means minimal number of edges, while in a
weighted graph it means minimal total edge weight. If several graph shortest paths exist, one
may choose any canonical representative, for example by a deterministic tie-breaking rule.

For an edge $e=\{u,v\}\in E$, define its length in the embedding $X$ by
\[
|X(e)| = \norm{\vect{x}_u-\vect{x}_v}.
\]
The embedding-induced geodesic distance along the fixed graph shortest path $\gamma_{ij}$
is then
\[
h_{ij}(X) = \sum_{e\in\gamma_{ij}} |X(e)|.
\]

This quantity is intrinsic to the drawn graph: it records the distance from $i$ to $j$
measured by traveling along graph edges, using the lengths induced by the embedding.

\section*{Geodesic KK Functional: Full All-Pairs Version}

The full geodesic KK functional uses all unordered vertex pairs and is defined by
\[
E_{\mathrm{geo}}(X)
=
\frac{1}{2}\sum_{1 \le i < j \le n}
k_{ij}\bigl(h_{ij}(X)-\ell_{ij}\bigr)^2,
\]
with the same target lengths and stiffnesses as in the classical model,
\[
\ell_{ij} = L_0 g_{ij},
\qquad
L_0 = \frac{L}{D},
\qquad
k_{ij} = \frac{K}{g_{ij}^2}.
\]

This is the direct analogue of the classical KK energy. It is the full model in the sense that
every chosen shortest path $\gamma_{ij}$ contributes a term, and each such term uses all
embedded edges belonging to that path through the quantity $h_{ij}(X)$.

This energy has a simple interpretation. The quantity $g_{ij}$ is the shortest-path distance in
the underlying graph metric. In the unweighted case, or when all graph edge weights are equal
to one, $g_{ij}$ is simply the number of edges in a shortest path from $i$ to $j$. In the
general weighted case, $g_{ij}$ is the sum of the prescribed graph edge lengths along a
shortest path. Accordingly, the target value $\ell_{ij}=L_0 g_{ij}$ says that the
embedding-induced geodesic distance along $\gamma_{ij}$ should be proportional to the
original graph metric, with proportionality scale $L_0$. The functional therefore promotes
drawings in which fixed graph shortest paths are realized by embedded polygonal chains whose
total lengths are approximately proportional to their graph distances.

Unlike the classical KK term, which only sees the chord $\norm{\vect{x}_i-\vect{x}_j}$, the
geodesic term sees the actual route through the graph. This is especially attractive when
the graph should be interpreted as a geometric network whose internal metric matters.

It is useful to note that
\[
h_{ij}(X)
=
\sum_{e\in\gamma_{ij}} \norm{\vect{x}_{u(e)}-\vect{x}_{v(e)}},
\]
so the functional is a stress model on sums of edge lengths along fixed shortest paths. Since
the paths $\gamma_{ij}$ are fixed, the derivatives can be assembled explicitly from path-edge
incidence information. In particular, every pair $(i,j)$ contributes only through the edges
contained in $\gamma_{ij}$, which makes implementation straightforward once the chosen
paths are stored.

\section*{Landmark-Sparse Geodesic KK Functional}

The all-pairs functional is expensive for large graphs because it uses all $\binom{n}{2}$
vertex pairs. A natural sparse alternative is to retain local interactions together with a small
set of long-range landmark constraints.

For each vertex $i$, let $N(i)$ denote the set of graph neighbors of $i$, and let $L(i)$ be a
set of landmark vertices associated with $i$. The landmark set should contain a small number
of vertices that are far from $i$ and well separated from one another. One convenient choice
is farthest-point sampling in the graph metric:
\begin{align*}
l_1(i)
&\in \operatorname*{arg\,max}_{v\in V} g_{iv}, \\
l_{t+1}(i)
&\in \operatorname*{arg\,max}_{v\in V}
\min\bigl\{g_{iv}, g_{l_1(i)v}, \dots, g_{l_t(i)v}\bigr\}.
\end{align*}
Then
\[
L(i) = \{l_1(i),\dots,l_M(i)\}
\]
for a prescribed landmark count $M$.

To avoid double counting, let
\[
\mathcal{S}
=
\bigl\{(i,j)\,:\, 1\le i<j\le n,\ \text{$j\in N(i)\cup L(i)$ or $i\in N(j)\cup L(j)$}\bigr\}.
\]
The landmark-sparse geodesic KK functional is
\[
E_{\mathrm{land}}(X)
=
\frac{1}{2}\sum_{(i,j)\in\mathcal{S}}
k_{ij}\bigl(h_{ij}(X)-\ell_{ij}\bigr)^2.
\]

This sparse model keeps two kinds of geometric information:
\begin{itemize}
\item local edge-scale information through the neighbors $N(i)$,
\item global shape information through the far-away landmark vertices $L(i)$.
\end{itemize}

The local terms encourage reasonable edge lengths and preserve nearby structure, while the
landmark terms prevent the layout from drifting globally. In practice, this can reduce both
memory use and optimization cost from quadratic scale to something much closer to linear
in the number of edges plus the number of landmark constraints.

\section*{Remarks}

The proposed geodesic KK functional may be viewed as replacing the ambient metric of the
drawing by the metric of the embedded graph itself. In this sense it is more intrinsic than the
classical KK model. The all-pairs version gives the clean conceptual formulation, while the
landmark version gives a practical large-scale approximation.

Both versions retain the familiar KK ingredients $g_{ij}$, $\ell_{ij}$, and $k_{ij}$; the only
change is the replacement of the naive Euclidean distance
\[
\norm{\vect{x}_i-\vect{x}_j}
\]
by the embedding-induced geodesic distance
\[
h_{ij}(X)=\sum_{e\in\gamma_{ij}} |X(e)|.
\]
This substitution aligns the energy with the internal geometry of the graph and therefore
better reflects the metric structure that the drawing is intended to represent.

\bibliographystyle{unsrt}
\bibliography{geodesic_kk_refs}

\end{document}
